\section{Positivités, espaces physiques et transport du spectre}
La géométrie positive et la théorie de jauge reçoivent des structures distinctes d'un même processus généalogique. La première réalise les poids et leurs incidences au moyen de mineurs ordonnés, de cellules et de formes canoniques. La seconde conserve l'incidence de jauge, l'holonomie, les représentations, les chemins et la mémoire au sein d'une mesure compatible. Sa forme de positivité pertinente pour la reconstruction physique est la positivité par réflexion. Le lien conceptuel acquiert un contenu mathématique lorsque les applications qui produisent chaque forme sont conservées.

Dans la partie géométrique, on a démontré que la subdivision d’une cellule
conserve sa forme canonique par annulation des pôles intérieurs.
Pour les énergies nodales, la réduction de Schur conserve une forme
effective et enregistre le détail. Dans la partie de jauge, la compatibilité du
raffinement doit conserver en outre le vide, l’espace des invariants,
le générateur et le témoin spectral. Telle est la chaîne développée
ci-dessous : échelle arithmético--géométrique, mesure commune, semi-groupe,
projection physique, raffinement et reconstruction continue.

La norme d’un vecteur, la positivité d’une matrice de moments et
l’existence d’un écart d’un Hamiltonien ont des domaines différents.
Une congruence à facteur inversible conserve la positivité et
l’inertie d’une forme ; un entrelacement isométrique du générateur, avec son
domaine, transporte sa réalisation sur l’image réduite de
l’inclusion. La borne uniforme et le témoin persistant fixent ensuite
l’écart limite. Le développement maintient donc visibles les
inclusions et les identités de semi-groupe qui justifient le passage à la
limite. L’échelle de l’écart provient de l’incidence aréale--volumétrique
et de son unité dimensionnelle déclarée.

Les développements suivants réunissent la construction projective et son développement ultérieur dépendant du couplage \(g\). Les premières comparaisons avec un paramétrage de Gibbs sont complétées ensuite au moyen de l'action, de la mesure et du transport explicites. L'ordre d'exposition conserve les démonstrations du processus et situe la formulation complète après ses dépendances. Les quatre réflexions correspondent aux quatre directions euclidiennes de la reconstruction. Le développement réuni ici a pour objet cette théorie de jauge et son écart spectral, avec le groupe compact, connexe et simple déclaré dans les théorèmes.

Le secteur supersymétrique planaire dans lequel les amplitudes sont étudiées au moyen d'amplituèdres conserve une question physique spécifique. Pour le relier à la présente réalisation, les opérateurs de supercharge, leurs domaines et la représentation correspondante sont requis. La dimension quatre de la coordinatisation euclidienne et la différence quatre du bloc modulaire demeurent des résultats déjà définis de leurs opérateurs respectifs.
